% LTeX: enabled=false
% ============================================================
% configuracoes/configuracoes-pdf.tex
%
% Configurações finais do PDF (hyperref, índice, hifenização).
% Carregado APÓS pacotes.tex (hyperref já foi carregado lá).
% ============================================================

% ── hyperref ─────────────────────────────────────────────────
\makeatletter
\hypersetup{%
    portuguese,
    hidelinks,          % sem cor e sem caixa; o negrito abaixo diferencia os links
    %    colorlinks  = true,
    breaklinks  = true,
    %    linkcolor   = \ifbool{darkmode}{DarkModeLink}{LightModeLink},
    %    citecolor   = \ifbool{darkmode}{DarkModeLink}{LightModeLink},
    %    filecolor   = \ifbool{darkmode}{DarkModeLink}{LightModeLink},
    %    urlcolor    = \ifbool{darkmode}{DarkModeLink}{LightModeLink},
    pdftitle    = {\@title},
    pdfauthor   = {\@author},
    pdfkeywords = {abnt, latex, abntex, abntex2, utfpr},
}
\makeatother

% ── Links em negrito (substituição de \hyperref colorido) ────
% Redefine o hook interno do hyperref para injetar \bfseries
% no início de cada link. Afeta \ref, \cite, \url, \href, etc.
\makeatletter
\let\UTFPR@hyper@linkstart\hyper@linkstart
\renewcommand{\hyper@linkstart}[2]{%
    \UTFPR@hyper@linkstart{#1}{#2}\bfseries
}
\makeatother

% ── Rótulos de \autoref ──────────────────────────────────────
\renewcommand{\algorithmautorefname}{Algoritmo}
\def\equationautorefname~#1\null{Equa\c{c}\~ao~(#1)\null}

% ── Índice remissivo ─────────────────────────────────────────
\makeindex

% ── Hifenização de palavras não reconhecidas ─────────────────
\hyphenation{%
    qua-dros-cha-ve
    Kat-sa-gge-los
}
